\documentclass[11pt]{article}

\usepackage[margin=1in]{geometry}
\usepackage[T1]{fontenc}
\usepackage[utf8]{inputenc}
\usepackage{authblk}

\title{Inferring Latent Engagement States in Decision-Making Under NMDA Receptor Modulation}
\author[1]{Javier Rodr\'iguez Mart\'inez}
\author[1]{Alexis Cervan}
\author[1]{Balma Serrano}
\author[2]{Alex Hyafil}
\author[2,3]{Gemma Huguet}
\author[1]{Jaime de la Rocha}
\affil[1]{Institut d'Investigacions Biom\`ediques August Pi i Sunyer (IDIBAPS)}
\affil[2]{Centre de Recerca Matem\`atica}
\affil[3]{Universitat Politècnica de Catalunya}

\date{}

\begin{document}

\maketitle

\begin{abstract}
Decision-making behavior is often analyzed as stationary, despite ongoing transitions between internal states such as engagement and disengagement that generate structured variability across trials. These history-dependent changes in behavioral strategy are not well captured by classical generalized linear models (GLMs), even when augmented with lapse terms, making latent-state approaches essential for interpreting performance and linking pharmacological perturbations to cognition.

In this project, we use hidden Markov models to infer latent engagement states from behavioral data collected in a two-alternative forced-choice task, a two-alternative delayed-choice task, and a multiple-choice delayed-response task. This framework allows us to describe behavior as a sequence of transitions between discrete internal states, each associated with distinct patterns of choice and performance. Rather than treating errors or variability as undifferentiated noise, the model provides an interpretable account of how internal state dynamics shape observable behavior over time.

Our current goal is to determine how pharmacological manipulation of NMDA receptors (NMDAr) affects these latent states and the transitions between them. By comparing behavioral data across drug conditions, we investigate whether NMDAr modulation changes the occupancy, stability, or switching dynamics of engagement-related states. In particular, we ask whether these manipulations promote prolonged disengagement, destabilize task-focused behavior, or induce distinct modes of responding that would remain hidden under conventional stationary analyses.

This approach offers a principled way to connect pharmacological interventions with internal cognitive dynamics, and may help clarify the contribution of NMDA-dependent mechanisms to adaptive decision-making. More broadly, it provides a quantitative framework for studying how neuromodulatory and pharmacological factors shape behavior through their effects on latent cognitive state structure.
\end{abstract}

\end{document}
